% !TeX root = ../main.tex
% 4. fejezet — Közös rendszeralapok
% A ftsrg/thesis-template-latex sablonba \include-dal illesztendő tartalomfájl.
% Hivatkozott források: Backend/HUNting.API/Program.cs, Data/HuntingContext.cs, Data/SeedData.cs,
%   Configuration/JwtOptions.cs, Services/UploadStorage.cs, Controllers/ApiControllerBase.cs,
%   Services/AdService.cs, SocietyService.cs, Controllers/* (végponttábla, feltöltési minta),
%   Migrations/*, appsettings.json, Properties/launchSettings.json,
%   Frontend/HUNting/src/main.tsx, App.tsx, routes/appRouter.tsx, layouts/AppLayout.tsx,
%   AuthLayout.tsx, lib/apiClient.ts, queryClient.ts, authStorage.ts, context/AuthContext.tsx,
%   components/layout/AppHeader.tsx, index.css

\chapter{Közös rendszeralapok}
\label{chap:alapok}

A HUNting funkciói -- a felhasználókezeléstől a térképes modulig -- közös infrastruktúrára
épülnek: ugyanazt az adatbázis-kontextust, ugyanazt a hitelesítési middleware-t és ugyanazt a
kliensoldali HTTP-réteget használják. Ez a fejezet ezt a funkciófüggetlen alapréteget mutatja be
egyszer, összefüggően, hogy a dolgozat gerincét adó
\ref{chap:felhasznalok}--\ref{chap:chatbot}.~funkciófejezetek ismétlés nélkül, kereszthivatkozással
építhessenek rá. A tárgyalás a rendszer rétegeit alulról felfelé követi: először az adattárolás
alapjait (adatbázis-kontextus, migrációk, kezdeti adatfeltöltés), majd a backend közös elemeit
(alkalmazás-összeállítás, rétegzés, JWT-infrastruktúra, visszatérő megvalósítási minták), végül a
frontend vázát (provider-fa, útvonalkezelés, API-kommunikáció, megjelenés) ismerteti.

\section{Adattárolási alapok}
\label{sec:alapok-adat}

\subsection{A \texttt{HuntingContext} és az adatmodell}

Az adatelérés az Entity Framework Core \emph{Code First} megközelítésére épül: az adatbázisséma
egyetlen hiteles forrása a C\# entitásmodell, amelyből a séma migrációkkal áll elő. A fejlesztői
környezetben az adatbázis egy \emph{VadászNet} nevű SQL Server LocalDB-példány, amelyet a
konfiguráció \texttt{DefaultConnection} kapcsolati karakterlánca jelöl ki.

A rendszer egyetlen adatbázis-kontextusa a \texttt{HuntingContext}
(\texttt{Data/HuntingContext.cs}), amely nem a sima \texttt{DbContext}, hanem az
\texttt{IdentityDbContext<ApplicationUser, IdentityRole<int>, int>} osztályból származik: így az
ASP.NET Core Identity felhasználó- és szereptáblái ugyanabban a sémában élnek, mint az üzleti
entitások, a teljes modell pedig egységesen egész számú elsődleges kulcsokat használ. A kontextus
húsz \texttt{DbSet} tulajdonságot tesz közzé, amelyek tartományok szerint csoportosíthatók:

\begin{itemize}[nosep]
	\item \emph{Felhasználók:} \texttt{Users} (az Identity alaptábláját fedi le);
	\item \emph{Társaságok:} \texttt{Societies}, \texttt{SocietyRoles}, \texttt{SocietyMembers},
	      \texttt{SocietyMembershipRequests}, \texttt{SocietyPosts};
	\item \emph{Hirdetések:} \texttt{Ads}, \texttt{AdCategories}, \texttt{AdImages},
	      \texttt{AdStatuses};
	\item \emph{Hírek:} \texttt{NewsFeed}, \texttt{NewsVisibility};
	\item \emph{Vadásznapló:} \texttt{HuntingLogs};
	\item \emph{Térkép:} \texttt{MapVisibilities}, \texttt{MapIconTypes}, \texttt{MapMarkers},
	      \texttt{MapPolygons}, \texttt{MapPolygonPoints}, \texttt{MapPaths},
	      \texttt{MapPathPoints}.
\end{itemize}

Ezekhez az Identity további hat infrastruktúra-táblája (szerepek, szerep-hozzárendelések,
claimek, külső bejelentkezések, tokenek) társul, így az adatbázis összesen 26 táblából áll. Az
Identity alapértelmezett \texttt{AspNet*} előtagú táblaneveit a kontextus
\texttt{OnModelCreating} metódusa beszédesebb nevekre (\texttt{Users}, \texttt{Roles},
\texttt{UserRoles}, \texttt{UserClaims}, \texttt{UserLogins}, \texttt{RoleClaims},
\texttt{UserTokens}) képezi le.

A modellkonfiguráció a Fluent API-ra és az entitásokon elhelyezett DataAnnotations
attribútumokra épül; a visszatérő megoldások a következők.
A státusz- és láthatóságjellegű \emph{kódtáblák} (\texttt{AdStatus}, \texttt{NewsVisibility},
\texttt{MapVisibility}) karakterlánc-kódokkal (például \texttt{ACTIVE}, \texttt{PUBLIC},
\texttt{MEMBERS}) kapcsolódnak a hivatkozó entitásokhoz: az idegen kulcs nem a numerikus
azonosítóra, hanem a kódoszlopra mutat (a \texttt{HasPrincipalKey} beállítással), így az
adatbázis-rekordok és az API-válaszok önmagukban is olvashatók. A kapcsolótáblák
(\texttt{SocietyMember}, \texttt{SocietyMembershipRequest}) \emph{összetett elsődleges kulcsot}
kapnak a társaság- és felhasználó-azonosító párosából, ami sémaszinten zárja ki a duplikált
tagságot és kérelmet. Az üzleti egyediséget \emph{egyedi indexek} garantálják (felhasználói és
társasági e-mail-cím, társasági szerepnév, hirdetéskategória-név, valamint a térképi pontoknál a
poligon-, illetve útvonal-azonosító és a \texttt{Seq} sorszám párosa). A létrehozási időbélyegek
(\texttt{CreatedAt}, \texttt{JoinedAt}) adatbázisoldali \texttt{sysdatetime()} alapértéket kapnak.
Végül a hirdetéseknél egy \emph{ellenőrző megszorítás} (\texttt{CK\_Ads\_SingleOwner}) mondja ki,
hogy minden hirdetésnek pontosan egy tulajdonosa van: vagy felhasználó, vagy társaság, de sosem
mindkettő vagy egyik sem. A teljes adatmodellt a~\ref{fig:er-diagram}.~ábra foglalja össze; az
egyes tartományok entitásait a megfelelő funkciófejezetek adatmodell-alfejezetei részletezik.

% TODO: E-01 ábra (figures/e-01-er-diagram.pdf) — teljes ER-diagram, Mermaid erDiagram kódból
% generálható a HuntingContext.cs alapján; egész oldalas, fekvő tájolás javasolt.
\begin{figure}[htbp]
	\centering
	\fbox{\parbox[c][9cm][c]{0.95\textwidth}{\centering E-01: A VadászNet adatbázis ER-diagramja
	(elkészítendő) -- Identity-mag + társasági, hirdetési, hír-, napló- és térképi tartomány}}
	\caption{A VadászNet adatbázis entitás--kapcsolat diagramja}
	\label{fig:er-diagram}
\end{figure}

\subsection{Törlési viselkedések}

A kapcsolatok törlési viselkedését a kontextus kapcsolatonként, kifejezetten rögzíti; ez a
relációs integritás szempontjából a modell egyik legfontosabb tervezési döntéssorozata.
A~\ref{tab:delete-behavior}.~táblázat a három alkalmazott viselkedés szerint csoportosítva
tekinti át a kapcsolatokat.

\begin{table}[htbp]
	\centering
	\caption{Törlési viselkedések a \texttt{HuntingContext} kapcsolataiban}
	\label{tab:delete-behavior}
	\begin{tabular}{|l|p{10.2cm}|}
		\hline
		\textbf{Viselkedés} & \textbf{Kapcsolatok} \\
		\hline
		\texttt{Cascade} &
		tagság és kérelem a társaság felől; tagság a felhasználó felől; társasági bejegyzés a
		társaság és a szerző felől; hirdetés a tulajdonos felhasználó, illetve társaság felől;
		hirdetéskép a hirdetés felől; hír a társaság és a szerző felől; naplóbejegyzés a
		felhasználó felől; térképi objektum a társaság felől; poligon- és útvonalpont a
		poligon, illetve útvonal felől \\
		\hline
		\texttt{Restrict} &
		tagság a társasági szerep felől; kérelem a kérelmező felhasználó felől; hirdetés a
		kategória és a státusz felől; hír a láthatósági kód felől; térképi objektum a
		létrehozó felhasználó, az ikontípus és a láthatósági kód felől \\
		\hline
		\texttt{SetNull} &
		a társaság \texttt{OwnerUserId} kulcsa a tulajdonos felhasználó törlésekor; a kérelem
		\texttt{ProcessedByUserId} kulcsa az elbíráló törlésekor \\
		\hline
	\end{tabular}
\end{table}

A mintázat következetes logikát tükröz. A \emph{tartalmazási} viszonyokban (a hirdetés képei, a
poligon pontjai, a társaság bejegyzései) a szülő törlése kaszkádban viszi magával a gyermekeket,
mert azok a szülő nélkül értelmezhetetlenek. A \emph{hivatkozási} viszonyokban -- ahol a
hivatkozott rekord megosztott kódtábla- vagy törzsadat-elem -- a \texttt{Restrict} akadályozza meg,
hogy egy még használatban lévő kategória, státusz vagy szerep eltűnjön a hivatkozó rekordok alól.
A \texttt{SetNull} ott jelenik meg, ahol a kapcsolat elvesztése nem érvényteleníti a rekordot:
a társaság a tulajdonosa nélkül is létező szervezet marad, az elbírált kérelem pedig az elbíráló
fiókjának megszűnése után is őrzi az állapotát. Megjegyzendő ugyanakkor egy aszimmetria: a térképi
objektumoknál a létrehozó felhasználóra mutató kulcs \texttt{Restrict} viselkedésű, vagyis a
felhasználó mindaddig nem törölhető, amíg térképi objektuma van -- miközben naplóbejegyzései és
hirdetései kaszkádban törlődnének. Ez a gyakorlatban azt jelenti, hogy az adminisztrátori
felhasználótörlés térképi tartalommal rendelkező fióknál adatbázishibára fut; a \texttt{Restrict}
viselkedésre a térképfejezet adatmodell-alfejezete (\ref{sec:map-adatmodell}.~alfejezet) is kitér.

\subsection{Migrációk}

A séma evolúcióját öt EF Core migráció írja le, amelyek egyben a fejlesztés idővonalát is
dokumentálják:

\begin{enumerate}[nosep]
	\item \texttt{InitialIdentity} (2025.~október) -- az Identity-alapú felhasználókezelés,
	      valamint a társasági, hirdetési, hír- és naplótartomány kezdeti sémája;
	\item \texttt{AddMapEntities} (2025.~november) -- a térképi entitások (láthatóságok,
	      ikontípusok, jelölők, poligonok, útvonalak és pontjaik) bevezetése;
	\item \texttt{SocietyAccountsAndPosts} (2025.~november) -- a társasági fiókok és a
	      társasági bejegyzések támogatása;
	\item \texttt{RenameIdentityTables} (2025.~december) -- az \texttt{AspNet*} táblanevek
	      lecserélése a beszédes nevekre;
	\item \texttt{AddDescriptionToMapPolygon} (2026.~március) -- leírásmező a térképi
	      poligonokhoz.
\end{enumerate}

A fejlesztői munkafolyamatban a migrációk a \texttt{dotnet ef migrations add} paranccsal
készülnek, alkalmazásukhoz azonban nincs szükség kézi \texttt{dotnet ef database update}
futtatásra: az alkalmazás indulásakor a kezdeti adatfeltöltés első lépése a
\texttt{Database.MigrateAsync()} hívás, amely a még hiányzó migrációkat automatikusan érvényesíti.
Ez a megoldás fejlesztői és demonstrációs környezetben kényelmes -- az adatbázis a backend első
indításával üzemkész --, éles, többpéldányos környezetben ugyanakkor a telepítéshez kötött,
ellenőrzött migrációfuttatás lenne indokolt.

\subsection{Kezdeti adatfeltöltés}

A \texttt{SeedData.SeedAsync} statikus metódus (\texttt{Data/SeedData.cs}) az alkalmazás minden
indulásakor lefut, és a migrációk érvényesítése után feltölti az adatbázist a működéshez és a
bemutatáshoz szükséges adatokkal. A feltöltés kétféle adatot érint. A \emph{törzsadatok} a rendszer
működéséhez szükségesek: a három alkalmazásszintű szerep (\emph{Administrator},
\emph{SocietyManager}, \emph{Hunter}), négy hirdetésstátusz, három-három hír- és térképláthatósági
kód, négy térképi ikontípus, négy társasági szerep (Vezető, Gazda, Vadőr, Tag) és tizenkét
hirdetéskategória. A \emph{demonstrációs adatok} a felületek bemutatását szolgálják: tizenegy
felhasználó (egy adminisztrátor és tíz vadász), nyolc vadásztársaság, tagsági kapcsolatok, hét
hirdetés képekkel, két hírbejegyzés és négy vadásznapló-bejegyzés.

A feltöltés kulcstulajdonsága az \emph{idempotencia}: minden rekord beszúrását egy természetes
kulcsra (e-mail-címre, névre, kódra vagy címre) szűrő létezés-ellenőrzés előzi meg, így az
ismételt indítás nem duplikál adatot, a hiányzó elemeket viszont pótolja.
A~\ref{lst:seed}.~kódrészlet ezt a mintát szemlélteti a társaságok feltöltésén.

\begin{lstlisting}[language={[Sharp]C}, caption={Idempotens adatfeltöltés a \texttt{SeedData} osztályban (kivonat)}, label={lst:seed}]
private static async Task EnsureSocietiesAsync(HuntingContext context)
{
    var societies = new[]
    {
        new Society { Name = "Bakonyi Vadasztarsasag",
                      Email = "kapcsolat@bakony-vt.example.hu", ... },
        // ... tovabbi het tarsasag ...
    };

    foreach (var society in societies)
    {
        if (!await context.Societies.AnyAsync(s => s.Email == society.Email))
        {
            context.Societies.Add(society);
        }
    }
}
\end{lstlisting}

A megoldás korlátaival együtt teljes a kép. A kezdő jelszavak (köztük az adminisztrátoré) a
forráskódban szerepelnek, ami kizárólag fejlesztői környezetben elfogadható. A
\emph{SocietyManager} szerep létrejön ugyan, de a feltöltött felhasználók egyike sem kapja meg: a
nyolc demonstrációs társaság közvetlenül, tulajdonos felhasználói fiók nélkül kerül az adatbázisba
(az \texttt{OwnerUserId} üresen marad, a kötelező \texttt{PasswordHash} mezőbe helykitöltő érték
kerül), társasági fiók így csak a regisztrációs folyamaton keresztül keletkezik
(\ref{chap:felhasznalok}.~fejezet). A hirdetésképek és a bejegyzésképek külső (Unsplash) URL-ekre
mutatnak, nem a lokális feltöltési könyvtárra.

\section{Backend-alapok}
\label{sec:alapok-backend}

\subsection{Az alkalmazás összeállítása}

A backend egyetlen belépési pontja a \texttt{Program.cs}, amely a .NET 8 minimál hoszting
modelljét követi: felső szintű utasításokkal, külön \texttt{Startup} osztály nélkül építi fel a
szolgáltatáskonténert és a HTTP-csővezetéket. A~\ref{lst:program-bootstrap}.~kódrészlet a fájl
szerkezetét mutatja a lényeges regisztrációkra és a middleware-sorrendre szűkítve.

\begin{lstlisting}[language={[Sharp]C}, caption={DI-regisztrációk és middleware-sorrend a \texttt{Program.cs}-ben (kivonat)}, label={lst:program-bootstrap}]
builder.Services.AddDbContext<HuntingContext>(opts =>
    opts.UseSqlServer(builder.Configuration
        .GetConnectionString("DefaultConnection")));

builder.Services.AddIdentityCore<ApplicationUser>(...)
    .AddRoles<IdentityRole<int>>()
    .AddEntityFrameworkStores<HuntingContext>()
    .AddSignInManager();

builder.Services.AddAuthentication(...).AddJwtBearer(...);
builder.Services.AddCors(...);

builder.Services.AddScoped<IAdService, AdService>();
builder.Services.AddScoped<ITokenService, JwtTokenService>();
// ... tovabbi service-regisztraciok ...
builder.Services.AddHttpClient<IOpenAiService, OpenAiService>();

var app = builder.Build();

if (app.Environment.IsDevelopment()) { app.UseSwagger(); app.UseSwaggerUI(); }
app.UseHttpsRedirection();
app.UseStaticFiles();      // wwwroot/uploads kiszolgalasa
app.UseCors("Frontend");
app.UseAuthentication();
app.UseAuthorization();
app.MapControllers();

await SeedData.SeedAsync(app.Services);
app.Run();
\end{lstlisting}

A szolgáltatásregisztrációk három csoportra oszthatók. Az \emph{infrastruktúra-szolgáltatások}
közé tartozik az adatbázis-kontextus, az Identity (jelszókezelés, szerepkezelés,
\texttt{SignInManager}) és a JWT-bearer hitelesítés. Az Identity jelszószabályzata szándékosan
enyhített: hat karakteres minimumhossz mellett sem kis- és nagybetűt, sem számjegyet, sem
speciális karaktert nem követel meg -- ez a demonstrációs célú használatot könnyíti, éles
üzemben szigorítandó. Az \emph{üzleti szolgáltatások} scoped életciklussal regisztrált
interfész--implementáció párok (\texttt{IAdService}, \texttt{IUserService},
\texttt{ISocietyService}, \texttt{IMapService}, \texttt{ITokenService}). A \emph{külső integrációt} az
\texttt{IOpenAiService} képviseli, amely típusos \texttt{HttpClient}-ként
(\texttt{AddHttpClient}) regisztrálódik, így a keretrendszer kezeli a HTTP-kapcsolatok
életciklusát (\ref{chap:chatbot}.~fejezet).

A kontrollerek JSON-szerializációja négy globális beállítást kap: a navigációs tulajdonságok
okozta körkörös hivatkozásokat a \texttt{ReferenceHandler.IgnoreCycles} vágja el, a \texttt{null}
értékű mezők kimaradnak a válaszokból, a szerializálási mélység legfeljebb 16 szint, a felsorolt
típusok pedig a \texttt{JsonStringEnumConverter} révén nevükkel (nem numerikus értékükkel)
utaznak -- így lesz a fióktípus a kliens számára \texttt{"Person"} vagy \texttt{"Society"}. A
\texttt{Frontend} nevű CORS-házirend a Vite fejlesztői kiszolgáló
(\texttt{http(s)://localhost:5173}) származási helyét engedélyezi tetszőleges fejléccel és
metódussal; a házirend a csővezetékben a hitelesítés \emph{előtt} fut le, így a böngészői
előkérések (preflight) tokenek nélkül is sikeresek. A middleware-sorrend további tagjai a
HTTPS-átirányítás, a statikus fájlok kiszolgálása (ez teszi elérhetővé a feltöltött képeket
tartalmazó \texttt{wwwroot/uploads} könyvtárat), majd a hitelesítés és a jogosultságkezelés. A
fejlesztői környezetben a Swagger (OpenAPI) felület is bekapcsolódik, amely a teljes
végpontkészletet dokumentálja és próbahívásokra is alkalmas
(\ref{fig:swagger}.~ábra).

% TODO: SC-17 képernyőkép (figures/sc-17-swagger.png) — Development módban, a végpontlistával.
\begin{figure}[htbp]
	\centering
	\fbox{\parbox[c][5.5cm][c]{0.85\textwidth}{\centering SC-17: A Swagger UI a HUNting.API
	végpontlistájával (elkészítendő)}}
	\caption{A generált Swagger-felület fejlesztői környezetben}
	\label{fig:swagger}
\end{figure}

\subsection{Rétegzés: kontrollerek és szolgáltatások}

A backend rétegzése a kontroller--szolgáltatás--adatelérés hármasra épül: a kontroller a
HTTP-szintű feladatokat (útvonalillesztés, státuszkódok, jogosultság-ellenőrzés) látja el, az
üzleti logika az interfész mögé rejtett szolgáltatásosztályokban él, az adatelérés pedig az EF
Core kontextuson keresztül történik. A~\ref{fig:komponens}.~ábra a kontrollerek, szolgáltatások
és a kontextus tényleges függőségi viszonyait mutatja.

% TODO: A-03 ábra (figures/a-03-komponensek.pdf) — komponensdiagram a Program.cs regisztrációi
% alapján: 8 controller → service-interfészek → HuntingContext; typed HttpClient → OpenAI.
\begin{figure}[htbp]
	\centering
	\fbox{\parbox[c][7cm][c]{0.9\textwidth}{\centering A-03: Backend-komponensdiagram
	(elkészítendő) -- kontrollerek, szolgáltatások, \texttt{HuntingContext}, külső OpenAI API}}
	\caption{A backend komponensei és függőségeik}
	\label{fig:komponens}
\end{figure}

Fontos azonban rögzíteni, hogy a rétegzés nem egységes. A szolgáltatásosztályok -- így az
\texttt{AdService} és a \texttt{SocietyService} -- közvetlenül a \texttt{HuntingContext}-tel
dolgoznak, külön repository-absztrakció nincs; a hírek és a vadásznapló kontrollere pedig
szolgáltatásréteg nélkül, közvetlenül a kontextusra épül. A kontrollerekben ismétlődő
segédlogika ugyanakkor két közös helyre szervezett: az \texttt{ApiControllerBase} ősosztály
adja a bejelentkezett felhasználó azonosítójának kiolvasását, valamint a szerep- és
fióktípus-ellenőrzéseket (\texttt{GetUserId}, \texttt{IsAdministrator},
\texttt{IsSocietyAccount}), a statikus \texttt{UploadStorage} segédosztály pedig a
képfeltöltések fájlrendszer-műveleteit végzi (lásd a közös megvalósítási mintáknál). Ahol egy
funkciófejezet a rétegzéstől eltér, ott azt külön jelzi.

\subsection{A JWT-infrastruktúra}

A hitelesítés állapotmentes, JSON Web Token alapú. A konfiguráció \texttt{Jwt} szakaszát a
\texttt{JwtOptions} beállításosztály (\texttt{Configuration/JwtOptions.cs}) képezi le, amely a
szimmetrikus aláírókulcsot, a kibocsátót (\texttt{HUNting.API}), a célközönséget
(\texttt{HUNting.Client}) és a token élettartamát (alapértelmezetten 60 perc) hordozza. A
bejövő kérések \texttt{Authorization} fejlécében érkező bearer-tokent a JWT-bearer middleware dolgozza
fel: ellenőrzi a kibocsátót, a célközönséget és az aláírást, a lejáratnál pedig mindössze
egyperces órakülönbséget tolerál. Sikeres validáció esetén a token claimjeiből előáll a kérés
identitása, amelyre a kontrollerek \texttt{[Authorize]} attribútumai és erőforrás-szintű
ellenőrzései épülnek. A tokenek \emph{kibocsátását} (a claimkészletet és az aláírást) a
felhasználókezelési fejezet (\ref{chap:felhasznalok}.~fejezet) mutatja be; ott esik szó a
megoldás korlátairól -- így a frissítő token és a tokenvisszavonás hiányáról -- is. Az
\texttt{appsettings.json}-ban tárolt aláírókulcs fejlesztői helykitöltő; éles környezetben
titokkezelő szolgáltatásból vagy környezeti változóból kellene érkeznie.

\subsection{Közös megvalósítási minták és a végpontkészlet}

Három megvalósítási minta a legtöbb kontrollerben visszatér, ezért itt kerül bemutatásra.

\paragraph{Validáció és hibakezelés.} Minden kontrollert az \texttt{[ApiController]} attribútum
jelöl, így a modellkötési és validációs hibákra a keretrendszer automatikusan 400-as választ ad.
A bemeneti szabályok első vonala az entitásokon elhelyezett DataAnnotations attribútumkészlet
(\texttt{[Required]}, \texttt{[MaxLength]} és társaik), amelyet az üzleti szabályok
kontrollerbeli, imperatív ellenőrzései egészítenek ki; a szolgáltatásrétegből érkező
\texttt{ArgumentException} jellegű hibákat a kontrollerek try/catch blokkjai fordítják 400-as
válaszra. Központi kivételkezelő middleware nincs: a nem kezelt kivétel a keretrendszer
alapértelmezett 500-as válaszát eredményezi.

\paragraph{Fájlfeltöltés.} A hirdetésképek, a társasági képek és a naplóképek feltöltése azonos
sémát követ: a végpontot a \texttt{[RequestSizeLimit(10\_485\_760)]} attribútum 10~MiB-ra
korlátozza; a kiterjesztés fehérlistán (\texttt{.jpg}, \texttt{.jpeg}, \texttt{.png},
\texttt{.webp}) ellenőrződik; a fájl ütközésmentes, GUID-alapú nevet kap, és a
\texttt{wwwroot/uploads} tartományspecifikus alkönyvtárába (\texttt{ads}, \texttt{societies},
\texttt{logs}) kerül; az adatbázisban csak a relatív URL tárolódik, amelyet a statikus
fájl-middleware szolgál ki; a rekord törlésekor a hozzá tartozó fájl is törlődik. A közös
lépéseket (kiterjesztés-ellenőrzés, könyvtár-előkészítés, mentés, törlés) a statikus
\texttt{UploadStorage} segédosztály (\texttt{Services/UploadStorage.cs}) végzi, így a
feltöltéskezelő kód nem ismétlődik kontrollerenként. A minta
korlátja, hogy a tartalom-ellenőrzés a kiterjesztésre szorítkozik (a fájl tényleges MIME-típusát
a szerver nem vizsgálja).

\paragraph{A végpontkészlet.} Az API nyolc kontrollere összesen 49 végpontot kínál;
a~\ref{tab:endpoint-osszes}.~táblázat ezek megoszlását összegzi. A részletes végponttáblák --
hozzáférési szabályokkal együtt -- az egyes funkciófejezetekben találhatók.

\begin{table}[htbp]
	\centering
	\caption{A HUNting.API kontrollerei és végpontjaik}
	\label{tab:endpoint-osszes}
	\begin{tabular}{|l|l|c|l|}
		\hline
		\textbf{Kontroller} & \textbf{Útvonalprefix} & \textbf{Végpont} & \textbf{Fejezet} \\
		\hline
		\texttt{AuthController} & \texttt{/api/Auth} & 4 & \ref{chap:felhasznalok}. \\
		\hline
		\texttt{UsersController} & \texttt{/api/Users} & 5 & \ref{chap:felhasznalok}. \\
		\hline
		\texttt{AdsController} & \texttt{/api/Ads} & 9 & \ref{chap:hirdetesek}. \\
		\hline
		\texttt{SocietiesController} & \texttt{/api/Societies} & 11 & \ref{chap:tarsasagok}. \\
		\hline
		\texttt{NewsController} & \texttt{/api/News} & 5 & \ref{chap:tarsasagok}. \\
		\hline
		\texttt{HuntingLogsController} & \texttt{/api/HuntingLogs} & 5 & \ref{chap:naplo}. \\
		\hline
		\texttt{MapController} & \texttt{/api/Map} & 9 & \ref{chap:terkep}. \\
		\hline
		\texttt{ChatController} & \texttt{/api/Chat} & 1 & \ref{chap:chatbot}. \\
		\hline
		\multicolumn{2}{|l|}{\textbf{Összesen}} & \textbf{49} & \\
		\hline
	\end{tabular}
\end{table}

\section{Frontend-alapok}
\label{sec:alapok-frontend}

\subsection{Belépési pont és provider-fa}

A kliensalkalmazás belépési pontja a \texttt{main.tsx}, amely a React 18 \texttt{createRoot}
API-jával, \texttt{StrictMode} burokban csatolja fel az \texttt{App} komponenst. Az \texttt{App}
(\texttt{App.tsx}) nem tartalmaz megjelenítési logikát: kizárólag az alkalmazást átfogó
szolgáltatók (providerek) fáját építi fel, amelyet a~\ref{fig:provider-fa}.~ábra szemléltet.
Legkívül a React Query \texttt{QueryClientProvider} áll, ezen belül az \texttt{AuthProvider}
biztosítja a hitelesítési állapotot, majd egy \texttt{Suspense} határ és a React Router
\texttt{RouterProvider} következik.

% TODO: FE-01 ábra (figures/fe-01-provider-fa.pdf) — fa-ábra: main → StrictMode → App →
% QueryClientProvider → AuthProvider → Suspense → RouterProvider → AppLayout/AuthLayout → oldalak.
\begin{figure}[htbp]
	\centering
	\fbox{\parbox[c][6cm][c]{0.85\textwidth}{\centering FE-01: A frontend provider- és
	layoutfája (elkészítendő)}}
	\caption{A kliensalkalmazás provider-hierarchiája és elrendezésrétegei}
	\label{fig:provider-fa}
\end{figure}

A szerverállapot kezelését a React Query végzi; a megosztott \texttt{queryClient} példány
(\texttt{lib/queryClient.ts}) három alapértelmezést rögzít: a lekérdezések eredménye egy percig
frissnek számít (\texttt{staleTime}), az ablakfókusz visszatérése nem vált ki újralekérést, a
sikertelen kérések pedig nem ismétlődnek automatikusan. A funkciómodulok konvenciója egységes: a
lekérdezések tartományspecifikus kulcsokkal (például \texttt{["ads", filters]}) regisztrálódnak,
a mutációk pedig sikeres lefutás után a \texttt{queryClient.invalidateQueries} hívással
érvénytelenítik az érintett kulcsok gyorsítótárát, így a felület automatikusan újratölti a friss
adatokat. Az \texttt{AuthProvider} és a mögötte álló \texttt{authStorage} modul működését -- a
token és a profil \texttt{localStorage}-beli tárolását, a munkamenet helyreállítását -- a
felhasználókezelési fejezet (\ref{chap:felhasznalok}.~fejezet) részletezi.

\subsection{Útvonalak és elrendezések}

Az útvonalkezelést a React Router \texttt{createBrowserRouter} hívása konfigurálja
(\texttt{routes/appRouter.tsx}), két elrendezés (layout) köré szervezve. Az \texttt{AppLayout}
a tartalmi oldalak kerete: rögzített fejlécet (\texttt{AppHeader}), legfeljebb \texttt{6xl}
szélességű tartalmi sávot, láblécet és -- minden oldalon elérhető, lebegő elemként -- a
chatbotot (\ref{chap:chatbot}.~fejezet) tartalmazza. Az \texttt{AuthLayout} a bejelentkezési és
regisztrációs oldal középre igazított kerete. A~\ref{tab:routes}.~táblázat a teljes
útvonaltáblát összegzi.

\begin{table}[htbp]
	\centering
	\caption{A kliensalkalmazás útvonalai}
	\label{tab:routes}
	\begin{tabular}{|l|l|l|}
		\hline
		\textbf{Útvonal} & \textbf{Oldal} & \textbf{Elrendezés} \\
		\hline
		\texttt{/} & \texttt{HomePage} & AppLayout \\
		\hline
		\texttt{/ads}, \texttt{/ads/new}, \texttt{/ads/:adId}, & hirdetésoldalak & AppLayout \\
		\texttt{/ads/:adId/edit} & & \\
		\hline
		\texttt{/societies}, \texttt{/societies/mine}, & társaságoldalak & AppLayout \\
		\texttt{/societies/manage}, \texttt{/societies/:societyId} & & \\
		\hline
		\texttt{/map} & \texttt{MapPage} & AppLayout \\
		\hline
		\texttt{/logs} & \texttt{HuntingLogsPage} & AppLayout \\
		\hline
		\texttt{/weather} & \texttt{WeatherPage} & AppLayout \\
		\hline
		\texttt{/profile} & \texttt{ProfilePage} & AppLayout \\
		\hline
		\texttt{/admin} & \texttt{AdminDashboardPage} & AppLayout \\
		\hline
		\texttt{/auth/login}, \texttt{/auth/register} & hitelesítési oldalak & AuthLayout \\
		\hline
		\texttt{*} & \texttt{NotFoundPage} & -- \\
		\hline
	\end{tabular}
\end{table}

Az ismeretlen útvonalakat egy globális, mindenre illeszkedő szabály és az \texttt{AppLayout}
\texttt{errorElement} beállítása egyaránt a \texttt{NotFoundPage} oldalra tereli. Két őszinte
megjegyzés kívánkozik ide. Egyrészt a router minden oldalkomponenst statikusan importál, azaz
\emph{nincs} kódszintű lusta betöltés (lazy loading): az \texttt{App} \texttt{Suspense} határa
jelenleg tartalék, egy későbbi kódfelbontás (code splitting) előkészítése, a teljes alkalmazás
egyetlen csomagban töltődik le. Másrészt az útvonaltábla nem tartalmaz központi útvonalőrt
(route guard): a védett oldalak hozzáférés-ellenőrzése oldalszinten, az egyes komponensekben
történik, ahogyan azt a felhasználókezelési fejezet bemutatja.

\subsection{API-kommunikáció}

A backenddel folytatott minden kommunikáció egyetlen megosztott axios-példányon, az
\texttt{apiClient}-en (\texttt{lib/apiClient.ts}) keresztül zajlik. Az alap-URL feloldása
háromlépcsős: ha a \texttt{VITE\_API\_BASE\_URL} környezeti változó be van állítva, az
érvényesül; fejlesztői módban az alapértelmezés a \texttt{https://localhost:5001/api} cím;
éles buildben a relatív \texttt{/api} útvonal. Megjegyzendő, hogy a backend fejlesztői profilja
valójában a 7003-as HTTPS-porton fut (\texttt{launchSettings.json}), így a fejlesztői
alapértelmezés önmagában nem talál célba -- a helyi fejlesztés a környezeti változó beállítását
igényli; ez a kódbázis ismert, javításra váró pontja. A példány viselkedésének magját a két
elfogó (interceptor) adja, amelyet a~\ref{lst:apiclient}.~kódrészlet mutat be.

\begin{lstlisting}[caption={Az \texttt{apiClient} kérés- és válasz-elfogója (kivonat)}, label={lst:apiclient}]
apiClient.interceptors.request.use((config) => {
  const token = getAccessToken();
  if (token) {
    config.headers.Authorization = `Bearer ${token}`;
  }
  return config;
});

apiClient.interceptors.response.use(
  (response) => response,
  (error) => {
    if (error.response?.status === 401) {
      purgeAuthStorage();
      window.location.href = "/auth/login";
    }
    if (error.response?.data?.message) {
      return Promise.reject(new Error(error.response.data.message));
    }
    // ... tovabbi hibauzenet-agak ...
  }
);
\end{lstlisting}

A kérés-elfogó a \texttt{localStorage}-ban tárolt hozzáférési tokent minden kimenő kérés
\texttt{Authorization} fejlécébe befűzi \texttt{Bearer} sémával, így az egyes funkciómodulok
API-hívásainak nem kell a hitelesítéssel foglalkozniuk. A válasz-elfogó két feladatot lát el:
a 401-es státuszú válaszokra globálisan reagál (törli a hitelesítési adatokat, és a
bejelentkezési oldalra irányít), minden más hibát pedig egységes \texttt{Error} objektummá
normalizál, a szerver \texttt{message} vagy \texttt{title} mezőjét részesítve előnyben -- a hívó
komponensek így mindig megjeleníthető hibaüzenetet kapnak. A modul harmadik közös eleme a
\texttt{resolveAssetUrl} segédfüggvény, amely a backend által visszaadott relatív képútvonalakat
(például \texttt{/uploads/ads/...}) teljes URL-lé egészíti ki az API gyökércíme alapján, az
abszolút (külső) URL-eket pedig változatlanul hagyja; ezt minden képet megjelenítő funkció
használja.

\subsection{Megjelenés és reszponzivitás}

A felület a Tailwind CSS 4-es verziójára épül, amelyben a téma nem külön konfigurációs fájlban,
hanem a globális stíluslapon (\texttt{index.css}), a \texttt{@theme} direktívával van definiálva.
A színvilág a vadászati tematikához igazodik: sötét erdőzöld alapszínek (háttér:
\texttt{\#111a13}, kártyák: \texttt{\#1f2a1f}) és arany-borostyán kiemelőszín
(\texttt{\#c3832f}) alkotják, a törzsszöveg az Inter, a címsorok a Playfair Display betűtípust
használják, a hátteret pedig rétegzett sugaras és lineáris színátmenet mélyíti. A reszponzivitás
mobil-első elvet követ, és a fejlécen figyelhető meg a legjobban: az \texttt{md} töréspont felett
vízszintes navigációs sáv és gombsor jelenik meg, alatta ezek elrejtőznek, és egy
hamburgermenü-gomb nyitja a lenyíló, érintésbarát navigációt. A fejléc rögzített, áttetsző,
elmosott hátterű sáv, a navigációs elemek pedig a bejelentkezett felhasználó fióktípusa és
szerepei szerint bővülnek (egyéni fióknál a „Társaságaim” és a „Naplók”, társasági fióknál a
„Társaságom”, adminisztrátornál az „Admin” menüpont). Megjegyzendő, hogy a felület magyar
feliratainak egy része jelenleg ékezet nélkül szerepel a forráskódban; ez ismert, javításra váró
szépséghiba.

% TODO: SC-01 képernyőkép (figures/sc-01-home-responsive.png) — főoldal desktop ÉS mobil
% nézetben egymás mellett; a mobil nézeten a hamburgermenü látszódjon.
\begin{figure}[htbp]
	\centering
	\fbox{\parbox[c][5.5cm][c]{0.85\textwidth}{\centering SC-01: A főoldal asztali és mobil
	nézetben (elkészítendő)}}
	\caption{A főoldal reszponzív megjelenése asztali és mobil nézetben}
	\label{fig:home-responsive}
\end{figure}

A fentiekben bemutatott alapréteg -- az egységes adatbázis-kontextus, a \texttt{Program.cs}-ben
összeállított backend-csővezeték és a provider-fára épülő kliensváz -- adja azt a közös nyelvet,
amelyre a következő fejezetek funkciónként építkeznek: mindegyikük az itt leírt mintákat
(JWT-validáció, \texttt{[Authorize]}, fájlfeltöltés, React Query-konvenciók,
\texttt{apiClient}) példányosítja a saját tartományára.
